% !PRINTED 170
\weberpage{170}
Es sei $F_1, F_2, \ldots, F_\mu$ eine nach II. bestimmte Basis des Systemes $\mathfrak{J}$, und $F$ irgend eine andere nicht constante Invariante von $S$. Dann lassen sich die Formen $A_1, A_2, \ldots, A_\mu$ so bebestimmen, dass % [sic]
\weq{2}{F = A_1F_1 + A_2F_2 + \cdots + A_\mu F_\mu}
wird. Wendet man auf diese identische Gleichung sämmtliche Substitutionen der Gruppe $S$ an, so bleiben nach Voraussetzung $F, F_1, F_2, \ldots F_\mu$ ungeändert, während $A_\varkappa$ in $A_\varkappa, A_\varkappa', A_\varkappa'' \ldots$ übergehen mag. Bildet man die Summe der so aus (2) abgeleiteten Gleichungen und setzt, wenn $m$ den Grad der Gruppe $S$ bedeutet,
\weq{3}{m\Phi_\varkappa = A_\varkappa + A_\varkappa' + A_\varkappa'' + \cdots,}
so folgt:
\weq{4}{F = \Phi_1F_1 + \Phi_2F_2 + \cdots + \Phi_\mu F_\mu.}
Die Functionen $\Phi_1, \Phi_2, \ldots \Phi_\mu$ sind aber nach \S.\,40 Invarianten von $S$, und damit ist nach III. die Endlichkeit des Systemes $\mathfrak{J}$ bewiesen.

Derselbe Schluss lässt sich auch auf die relativen Invarianten anwenden, wie folgt:

Wir bezeichnen mit $F(x)$ das ganze System der Functionen, die den Bedingungen \S.\,40, (3):
\[
F[A(x)] = \alpha F(x), \quad F[B(x)] = \beta F(x) \ldots
\]
bei feststehenden Factoren $\alpha, \beta, \ldots$, die, wie wir gesehen haben, Einheitswurzeln sind, genügen.

Nach dem \gesperrt{Hilbert'schen} Satze lässt sich ein speci\"elles System solcher Functionen $F_1, F_2, \ldots, F_m$ derart auswählen, dass man
\weq{5}{F = A_1F_1 + A_2F_2 + \cdots + A_mF_m}
setzen kann, worin $A_1, A_2, \ldots, A_m$ Formen der Variablen $(x)$ sind. Behandelt man diese Formel so wie die Formel (2), indem man die Substitutionen der Gruppe $S$ darauf anwendet und dann die Summe bildet, so erhält man, entsprechend der Formel (4):
\[
F = \Phi_1F_1 + \Phi_2F_2 + \cdots + \Phi_mF_m,
\]
worin die Co\"efficienten $\Phi_1, \Phi_2, \ldots, \Phi_m$ \gesperrt{absolute Invarianten} sind.

Zur Vervollständigung ist noch hinzuzufügen, dass inhomogene Functionen der Variablen nur dann Invarianten sein können, wenn ihre einzelnen homogenen Bestandtheile Invarianten sind.
